% Generated draft: requires mathematical and source review.
\documentclass[11pt]{article}
\usepackage[margin=25mm]{geometry}
\usepackage{fontspec}
\setmainfont{Latin Modern Roman}
\usepackage{amsmath,amssymb,mathtools,unicode-math}
\setmathfont{Latin Modern Math}
\usepackage{enumitem,needspace,xurl,hyperref}
\setlength{\parindent}{0pt}
\setlength{\parskip}{5pt}
\setlength{\emergencystretch}{3em}
\newcommand{\N}{\mathbb{N}}
\newcommand{\Z}{\mathbb{Z}}
\newcommand{\Q}{\mathbb{Q}}
\newcommand{\R}{\mathbb{R}}
\newcommand{\C}{\mathbb{C}}
\newcommand{\F}{\mathbb{F}}
\newcommand{\A}{\mathbb{A}}
\newcommand{\PP}{\mathbb{P}}
\newcommand{\E}{\mathbb{E}}
\newcommand{\eps}{\varepsilon}
\newcommand{\dd}{\,\mathrm{d}}
\begin{document}
\section*{UnsolvedMath ID: 30000052}
\textbf{Source problem code:} catalogue entry
\subsection{Definitions and mathematical statement}
Let \(d,n\in\N\) and \(p\in\{2,4,6,\ldots\}\). Choose \(Z,\Delta\) independently and uniformly from \([0,1)^d\), and set
\[
X_j=\{jZ+\Delta\}\quad(0\leq j<n),\qquad
M_n^{Z,\Delta}=(X_0,\ldots,X_{n-1}),
\]
where braces denote coordinatewise fractional part. Separately, let \(U_0,\ldots,U_{n-1}\) be independent uniform points in \([0,1)^d\), and put \(M_n=(U_0,\ldots,U_{n-1})\). This notation treats both configurations as \(n\) sampled points, with multiplicity; coincidences have probability zero for these distributions.

The source uses \(\operatorname{disc}_p\) without reproducing its definition. Under the usual anchored-box \(L^p\) convention, for \(P=(x_0,\ldots,x_{n-1})\) define
\[
D_P(t)=\frac1n\sum_{j=0}^{n-1}\mathbf1_{[0,t_1)\times\cdots\times[0,t_d)}(x_j)-\prod_{r=1}^d t_r,
\qquad
\operatorname{disc}_p(P)^p=\int_{[0,1)^d}|D_P(t)|^p\,\dd t.
\]
On this convention, the cleaned source question is whether, for every \(d,n\in\N\) and every even positive integer \(p\geq2\),
\[
\E_{Z,\Delta}\bigl[\operatorname{disc}_p(M_n^{Z,\Delta})^p\bigr]
\leq
\E\bigl[\operatorname{disc}_p(M_n)^p\bigr].
\]
The anchored-box convention must be verified in the report; alternative discrepancy definitions are not automatically equivalent. Replacing the normalized counting error by its unnormalized version multiplies both sides by \(n^p\), provided the same convention is used on both sides.

The original short prose speaks of expected \(L^p\) discrepancy, while the cleaned, reportedly source-verified statement explicitly compares expectations of its \(p\)-th power. The latter is retained here: an inequality for the expected norm is a different assertion and must not be substituted. The supplied literature review checked 2026-08-21 reports that the checked shifted-lattice literature did not prove this displayed comparison and labels it open.
\subsection{Short English statement}
For even \(p\), compare the expected \(p\)-th power of discrepancy for one uniformly randomized arithmetic progression on the torus with that for \(n\) independent uniform points. The exact discrepancy convention still needs confirmation from the report.
\subsection{Sources}
\begin{itemize}
\item Local UnsolvedMath registry, permanent ID 30000052, catalogue entry, \emph{Expected Discrepancy of Generalized Lattice Point Sets}.
\item \emph{Discrepancy Theory and Its Applications} (2004), Oberwolfach Reports, pp. 673--722, DOI \url{https://doi.org/10.4171/owr/2004/13}. The literature field also supplies \url{https://doi.org/10.4171/OWR/2004/13}.
\item Supplied reference labeled ``Shifted lattice rules'': \url{https://doi.org/10.1016/j.cam.2009.06.002}.
\item Supplied statement-verification field reports checking the official report; supplied literature triage is dated 2026-08-21. Neither check is independently repeated here.
\end{itemize}
\end{document}
